\documentclass[11pt]{amsart}
\usepackage[margin=1.2in]{geometry}
\usepackage{amsmath,amssymb,amsthm}
\usepackage{booktabs,longtable,array}
\usepackage[colorlinks=true,linkcolor=blue,citecolor=blue,urlcolor=blue]{hyperref}

\newtheorem{theorem}{Theorem}[section]
\newtheorem{proposition}[theorem]{Proposition}
\newtheorem{lemma}[theorem]{Lemma}
\newtheorem{corollary}[theorem]{Corollary}
\newtheorem{conjecture}[theorem]{Conjecture}
\theoremstyle{definition}
\newtheorem{definition}[theorem]{Definition}
\newtheorem{construction}[theorem]{Construction}
\newtheorem{example}[theorem]{Example}
\theoremstyle{remark}
\newtheorem{remark}[theorem]{Remark}
\newtheorem{boundary}[theorem]{Claim Boundary}

\newcommand{\id}{\mathrm{id}}
\newcommand{\PSh}{\mathrm{PSh}}
\newcommand{\Sh}{\mathrm{Sh}}
\newcommand{\Set}{\mathbf{Set}}
\newcommand{\Hom}{\mathrm{Hom}}
\newcommand{\Ob}{\mathrm{Ob}}
\newcommand{\Noe}{\mathrm{Noetica}}
\newcommand{\Jleg}{J_{\mathrm{legacy}}}
\newcommand{\Eamb}{\mathcal{E}_{\mathrm{amb}}}
\newcommand{\Eres}{\mathcal{E}_{\mathrm{res}}}
\newcommand{\Eleg}{\mathcal{E}_{\mathrm{legacy}}}
\newcommand{\cC}{\mathcal{C}}

\begin{document}

\title[The marked and modal Noetic topos]{Noetic Topos 2.0: a marked and modal successor to the conservative Noetic site, with a complete finite-context theory}
\author{Little Henderson}
\thanks{Research paper v0.2, 2026-08-07 (supersedes white paper v0.1). Artifact status: \textsc{draft / gate 1 research candidate / ring 4 / not licensed / held}. Source basis: \texttt{NOETIC\_TOPOS\_2\_SPEC\_v0.1.md} (SHA-256 \texttt{156A5D8C5320F36E81E3B46473BA8A038B35672DCD28839A404D96A3DB1A763A}); theorem ledger v0.1 (\texttt{52E00065\ldots}); proof-obligation ledger v0.1 (\texttt{FD567E38\ldots}); \texttt{FINITE\_MODEL\_CERTIFICATE\_v0.1.json}; \texttt{MODALITY\_COMPOSITION\_PROOF\_v0.1.md} and certificate; publication delta Paper~2 at base commit \texttt{a1cd990109e7561dfd03949bf0e1653461c7ae44}; implementation \texttt{src/noetic\_topos2/}. Every theorem below carries the proof status recorded in the ledger; drafting this paper promotes no claim, and independent category-theory review of every result remains open.}
\date{August 7, 2026}

\begin{abstract}
The conservative Noetic Topos construction contains a valid conditional topos-existence argument, but its stated ACBE and RPM cover clauses do not retain the separate information later interpretations assigned to them. ACBE reflexivity places an identity in the corresponding sieve; the RPM counit, split by its section, forces an identity by sieve closure. Each clause alone contributes only maximal sieves, and the combined topology reduces to the jointly-surjective coverage already present in the source. A successor cannot obtain distinct modalities by relabeling those generators.

This paper develops the Noetic Topos 2.0 finite reference candidate as a conservative repair, in full. It augments the legacy site with a four-atom context lattice $L=\mathcal{P}(\{0,1\}\times\{0,1\})$, equips that factor with two transverse coordinate-partition coverages, and keeps immutable, hash-chained cover provenance outside the ordinary topology. On $\PSh(L,\Set)$ we prove: both generated topologies are subcanonical; the associated-sheaf reflectors are distinct and non-identity, witnessed by explicit finite presheaves; their join is the topology generated by the four-singleton partition; an explicit product-over-blocks construction realizes each associated-sheaf reflector, verified directly against the universal property; and the two reflectors commute up to unit-compatible canonical natural isomorphism, with either composite naturally isomorphic to sheafification for the joined topology. The last result is corrective as well as constructive: the finite modalities are distinct but \emph{commuting}, and no artifact licenses a noncommutation claim. An exhaustive finite certificate (16 contexts, 81 inclusion arrows, 298 sieves, 168 covering sieves, 18 executable fixtures, 3{,}012 restriction squares) corroborates every finite theorem. The ambient candidate $\mathrm{NoeTop}_2=\PSh(\cC\times L)$ carries three modal reflectors and a resolved subtopos; the full research object packages these with an external mark record. The product-site lift, legacy round-trip comparison, cross-context coupling, proof-assistant formalization of the sheaf layer, and independent review remain open, and are stated as explicit obligations with acceptance criteria.
\end{abstract}

\maketitle
\tableofcontents

\section{Introduction}

AI-assessment work in the Cognitive Architecture System (CAS) program increasingly relies on heterogeneous evidence gathered under different prompts, tasks, roles, and deployment contexts. The recurring difficulty is structural rather than numerical: local observations need not compose coherently, and a discipline is needed for tracking which local views count as jointly informative. Category theory addresses precisely this: a \emph{site} records which families of maps count as covers, and the associated sheaf category enforces compatibility constraints on data collected over those covers. The conservative Noetic Topos paper introduced such a site for the registered category $\Noe$ and proved a conditional topos-existence theorem.

This paper is about what that construction could not do, why, and what the smallest honest repair looks like.

\subsection{The problem}
The conservative construction admits covering sieves from three clauses: jointly-surjective observational families, designated ACBE descent maps, and the RPM counit. The intent was that the ACBE and RPM clauses would later support distinct modal operators---separate ``lenses'' on evidence. Section~\ref{sec:collapse} shows this intent is unrealizable from those generators: both clauses generate only maximal sieves, so the topology they define collapses to the observational coverage, and ordinary sieve generation \emph{forgets} the ACBE/RPM provenance entirely. The collapse is not a defect of the topos machinery; it is a theorem about which information a Grothendieck topology can retain.

\subsection{The contribution}
The successor candidate makes three moves, each developed fully below.

\begin{enumerate}
\item \textbf{Change the generators} (Sections~\ref{sec:factor}--\ref{sec:theorems}). A finite context lattice $L$ with two transverse partition coverages replaces identity-containing clauses with proper-piece covers. We prove the generated topologies are genuine, subcanonical, and support distinct non-identity sheafification modalities, and we compute their join and composition laws exactly.
\item \textbf{Move provenance out of the topology} (Section~\ref{sec:marks}). Since no topology can remember why a cover was admitted, the candidate stores that information in immutable, hash-chained \emph{marked cover} records that travel in parallel with the unmarked families the topology consumes.
\item \textbf{Name the ambient package honestly} (Sections~\ref{sec:ambient} and \ref{sec:compat}). The successor object is the ambient presheaf topos with three modal reflectors and a resolved subtopos---not the resolved subtopos alone, on which the new modalities provably restrict to identity.
\end{enumerate}

\subsection{Evidence discipline}
Every theorem in this paper carries the proof status recorded in the registered theorem ledger: all finite-context theorems are \textsc{informal proof present} with exhaustive executable corroboration; the backward-compatibility statement is a \textsc{conjecture} with \textsc{proof required}; nothing is peer-reviewed or formalized beyond the bounded Lean scope stated in Section~\ref{sec:certificate}. The ACBE and RPM coordinate names are formal marks with no empirical interpretation; no theorem here establishes that either coordinate detects a cognitive, causal, safety, or alignment property.

\section{Preliminaries}\label{sec:prelim}

We recall the standard notions; see \cite{maclane,stacks00ZC}.

\begin{definition}\label{def:sieve}
Let $\cC$ be a category and $C\in\Ob(\cC)$.
\begin{enumerate}
\item A \emph{sieve} $S$ on $C$ is a collection of morphisms with codomain $C$ closed under precomposition: if $f:D\to C$ lies in $S$ and $h:E\to D$ is any morphism, then $f\circ h\in S$.
\item If $g:D\to C$ is a morphism and $S$ is a sieve on $C$, the \emph{pullback sieve} $g^{*}S$ on $D$ is $\{h:E\to D \mid g\circ h\in S\}$.
\item The \emph{maximal sieve} $m_C$ on $C$ is the class of all morphisms with codomain $C$.
\end{enumerate}
\end{definition}

\begin{definition}\label{def:topology}
A \emph{Grothendieck topology} $J$ on $\cC$ assigns to each object $C$ a collection $J(C)$ of sieves on $C$, called covering sieves, such that: (J1) $m_C\in J(C)$; (J2) if $S\in J(C)$ and $g:D\to C$, then $g^{*}S\in J(D)$; (J3) if $S\in J(C)$ and $T$ is a sieve on $C$ with $f^{*}T\in J(D)$ for every $f:D\to C$ in $S$, then $T\in J(C)$. The pair $(\cC,J)$ is a \emph{site}.
\end{definition}

\begin{definition}\label{def:sheaf}
Let $(\cC,J)$ be a site. A presheaf $F:\cC^{\mathrm{op}}\to\Set$ is a \emph{sheaf} if for every object $C$ and covering sieve $S\in J(C)$ the canonical map
\[
F(C)\longrightarrow \mathrm{Nat}(S,F)
\]
is a bijection, where $S$ is regarded as a subfunctor of $yC=\Hom_{\cC}(-,C)$. A topology is \emph{subcanonical} when every representable presheaf is a sheaf.
\end{definition}

\begin{remark}\label{rem:background}
If $(\cC,J)$ is a small site, then $\Sh(\cC,J)$ is a Grothendieck topos \cite[Chap.~VII]{maclane}. Sheafification $a:\PSh(\cC)\to\Sh(\cC,J)$ is left adjoint to inclusion $i$ and is left exact \cite{stacks00W1}. The composite $M=ia$ is a left-exact idempotent monad on $\PSh(\cC)$; we call such composites \emph{modalities} and say a modality is \emph{trivial} when it is naturally isomorphic to the identity. A coverage basis (pretopology) generates a topology, and a presheaf is a sheaf for the generated topology if and only if it satisfies the gluing condition for the basis families \cite{stacks00ZC}.
\end{remark}

\begin{remark}[Elementary but decisive facts]\label{rem:facts}
Two observations drive Section~\ref{sec:collapse}. (i) A sieve containing $\id_C$ is maximal: for any $f:D\to C$, $f=\id_C\circ f$ lies in the sieve by precomposition closure. (ii) If a sieve contains a split epimorphism $\varepsilon$ with section $\eta$, then it contains $\varepsilon\circ\eta=\id$, hence is maximal.
\end{remark}

\section{The conservative baseline and its information loss}\label{sec:collapse}

Let $\cC$ denote the registered legacy category $\Noe$: a small concrete category of finite sets indexed by the forty distinguished CAS indices, with admissibility-relation generating maps and pullback closure. Let $\Jleg$ denote the coverage defined in the conservative source, generated by three clauses: (O) jointly-surjective observational families; (A) designated ACBE descent maps, including the reflexive descent map at equal contexts; (R) the RPM counit $\varepsilon_X$, where the RPM data supply
\[
\eta_X : X \to R(X), \qquad \varepsilon_X : R(X)\to X, \qquad \varepsilon_X\circ\eta_X=\id_X ,
\]
a section-retraction pair, not an adjunction.

\begin{proposition}[Legacy collapse]\label{prop:collapse}
Under the registered axioms, every sieve generated by clause \textup{(A)} alone or clause \textup{(R)} alone is maximal. Consequently
\[
\Jleg \;=\; J_{\text{jointly-surjective}},
\]
and the generated topology retains no ACBE or RPM provenance.
\end{proposition}

\begin{proof}
For (A): the reflexive descent map at equal contexts is $\id_X$; by Remark~\ref{rem:facts}(i) the generated sieve is maximal. For (R): a sieve containing $\varepsilon_X$ contains $\varepsilon_X\circ\eta_X=\id_X$ by precomposition closure (Remark~\ref{rem:facts}(ii)); maximal again. Maximal sieves belong to every Grothendieck topology by (J1), so neither clause adds a nonmaximal cover, and the generated topology equals the one generated by clause (O) alone.
\end{proof}

\begin{remark}
Proposition~\ref{prop:collapse} does not invalidate the conditional topos-existence theorem of the source: $\Sh(\cC,\Jleg)$ is still a Grothendieck topos when $\cC$ is small. What collapses is the \emph{intended interpretive structure}: any modality built from the (A)- or (R)-generated topology is trivial, because those topologies contain only what joint surjectivity already gives. The v7.4 corrigendum therefore permits the topos statement while retracting the inference that the legacy clauses generate distinct nontrivial modalities. The registered finite certificate independently reconfirms both collapse mechanisms as executable checks (\texttt{legacy\_acbe\_identity\_forces\_maximal\_sieve}, \texttt{legacy\_rpm\_split\_epi\_forces\_identity}).
\end{remark}

\begin{proposition}[Non-derivability of legacy subcanonicity]\label{prop:countermodel}
Subcanonicity of $(\cC,\Jleg)$ is not derivable from the registered axioms.
\end{proposition}

\begin{proof}
It suffices to exhibit a model of the axioms whose jointly-surjective coverage is not subcanonical. Take the category of sets of size at most two with injections as morphisms. The two point inclusions $\{a\}\hookrightarrow\{a,b\}$, $\{b\}\hookrightarrow\{a,b\}$ are jointly surjective, hence generate a covering sieve. Consider the representable $y\{*\}$ evaluated along this cover: a matching family assigns to each point inclusion an injection into $\{*\}$ (there is exactly one from each singleton), and the compatibility condition on the empty overlap is vacuous, so the matching-family set is a singleton---but so is checked against \emph{gluings}: a global section would be an injection $\{a,b\}\to\{*\}$, and none exists. The gluing map is therefore not surjective onto matching families, and $y\{*\}$ is not a sheaf.
\end{proof}

\begin{boundary}
Proposition~\ref{prop:countermodel} establishes non-derivability, not universal non-subcanonicity of every possible $\Noe$ model. Until an exact revised site is proved subcanonical, the distinguished legacy object is represented in the sheaf topos by $a_{\mathrm{legacy}}(yX)$, never by silently treating $yX$ as a sheaf \cite{stacks00WO}.
\end{boundary}

\section{The finite marked context factor}\label{sec:factor}

The successor must change the cover generators. The repair introduces a separate finite \emph{context} dimension on which nontrivial coverage is possible, leaving the legacy category untouched.

\begin{definition}[Context lattice]\label{def:lattice}
Let
\[
\Omega=\{(0,0),(0,1),(1,0),(1,1)\},\qquad L=\mathcal{P}(\Omega),\qquad (U\to V \iff U\subseteq V).
\]
$L$ is a sixteen-object poset with $81$ inclusion arrows. The first coordinate of an atom is the formal ACBE mark, the second the formal RPM mark. Pullbacks in $L$ are intersections: the pullback of $U\to V$ along $W\to V$ is $U\cap W$.
\end{definition}

\begin{boundary}
The coordinate names ACBE and RPM are formal labels preserved for controlled migration vocabulary. They carry no empirical interpretation until a separate operationalization and validation record exists. No theorem in this paper attaches meaning to them beyond their combinatorics.
\end{boundary}

\begin{definition}[Transverse partitions and generating families]\label{def:covers}
Define the two transverse partitions of $\Omega$
\[
A_i=\{(i,0),(i,1)\}, \qquad R_j=\{(0,j),(1,j)\}, \qquad i,j\in\{0,1\},
\]
and, for $P\in\{A,R\}$ with blocks $B$, the generating family at a nonempty context $V$
\[
K_P(V)=\{\,V\cap B \to V \mid B \text{ a block of } P\,\},
\]
retaining an empty piece when an intersection is empty. At the empty context the generator is the null (empty) family, which covers it. Let $J_A$ and $J_R$ be the generated topologies, and let $Q$ denote the four-block singleton partition with the same null-cover convention, generating $J_{\mathrm{atom}}$.
\end{definition}

The transversality identity that powers everything below is
\begin{equation}\label{eq:transverse}
A_i\cap R_j=\{(i,j)\}.
\end{equation}

\begin{lemma}[Coverage basis]\label{lem:basis}
For each $P\in\{A,R,Q\}$, the families $K_P$ together with identity families and the null cover form a pullback-stable coverage basis: pullback along $U\subseteq V$ carries $K_P(V)$ to $K_P(U)$ (as an indexed family; to identity arrows on the empty object when $U=\varnothing$), and refining a member by the same partition regenerates the same sieve.
\end{lemma}

\begin{proof}
Pullback of $V\cap B\to V$ along $U\subseteq V$ is $(V\cap B)\cap U=U\cap B\to U$. If $U\neq\varnothing$ this is exactly the indexed family $K_P(U)$; if $U=\varnothing$ each pullback is $\id_\varnothing$, and the registered null cover covers $\varnothing$ as well. For composition: refining a member $V\cap B$ by the same partition replaces it with $\{V\cap B\cap B' \mid B'\in P\}$; block disjointness makes every intersection with $B'\neq B$ empty, so the only nonempty piece is $V\cap B$ itself, and the composite family generates the same sieve. Identity covers supply the remaining basis axiom.
\end{proof}

\begin{theorem}[NT2-1: generated topologies; \textnormal{ledger \textsc{thm-nt2-marked-coverage}}]\label{thm:nt2-1}
$K_A$ and $K_R$ generate Grothendieck topologies $J_A$ and $J_R$ on $L$.
\end{theorem}

\begin{proof}
Immediate from Lemma~\ref{lem:basis} and the standard passage from a coverage basis to covering sieves \cite{stacks00ZC}: the generated collection satisfies (J1) by identity covers, (J2) by pullback stability, and (J3) by closure under cover composition.
\end{proof}

\begin{proposition}[NT2 global covers are proper; \textnormal{ledger \textsc{thm-nt2-nontrivial-global-covers}}]\label{prop:proper}
At the global context $\Omega$, both pieces of $K_A(\Omega)=\{A_0\to\Omega,\ A_1\to\Omega\}$ and of $K_R(\Omega)$ are proper subcontexts; the generated sieves are nonmaximal.
\end{proposition}

\begin{proof}
$A_i$ and $R_j$ are two-element subsets of the four-element $\Omega$, hence proper. A sieve on $\Omega$ generated by arrows from proper subobjects of a poset cannot contain $\id_\Omega$: every arrow in the generated sieve factors through some $A_i$ (respectively $R_j$), and $\id_\Omega$ factors through no proper subset. By Remark~\ref{rem:facts}(i) the sieve is nonmaximal.
\end{proof}

This is the exact point at which the successor escapes Proposition~\ref{prop:collapse}: unlike the legacy clauses, the new generators contain no identity and no split epimorphism onto their target.

\section{Marks as parallel audit data}\label{sec:marks}

Proposition~\ref{prop:collapse} shows a topology cannot remember why a family was admitted; sieve generation is a closure operator that erases provenance. The candidate therefore maintains two representations with distinct consumers.

\begin{construction}[Marked covers]\label{con:marks}
An \emph{unmarked family} is the exact list of arrows consumed by topology generation. A \emph{marked cover} is the immutable record
\[
\bigl(\text{derivation\_kind},\ \text{target},\ \text{exact\_member\_arrows},\ \text{mark\_record},\ \text{parent\_digests}\bigr),
\]
with $\text{derivation\_kind}\in\{\textsc{generator},\textsc{pullback},\textsc{common refinement}\}$ and mark record
\[
\bigl(\text{cover\_kind},\ \text{source\_authority\_ref},\ \text{assumptions},\ \text{evidence\_refs},\ \text{review\_state},\ \text{provenance\_chain}\bigr).
\]
Canonical JSON with SHA-256 binds target, members, mark, and the sorted set of parent-cover digests. Pullback along $U\to V$ creates a new record whose provenance chain contains its parent digest; the ACBE/RPM common refinement is a derivation node with both generating-cover digests as parents.
\end{construction}

The integrity rules are fail-closed. Unknown or duplicate JSON keys and uncontrolled review states are rejected. Deserialization of a derived record requires a parent resolver: every parent digest is rehashed, the complete acyclic parent derivation is checked, and the child members and mark are recomputed; a common refinement without exactly one ACBE and one RPM generating parent, or a pullback without its exact parent, is rejected rather than treated as evidence.

\begin{example}[Why marks are necessary even at singletons]\label{ex:singleton}
At a singleton context $\{(i,j)\}$, the bare families coincide:
\[
K_A(\{(i,j)\})=\{\{(i,j)\}\to\{(i,j)\},\ \varnothing\to\{(i,j)\}\}=K_R(\{(i,j)\})
\]
as unmarked arrow lists. The generated sieves are equal, so the topology cannot distinguish them---but the marked records carry different cover kinds and different digests, and the certificate confirms both \texttt{singleton\_unmarked\_families\_coincide} and \texttt{singleton\_marked\_covers\_distinct}. Provenance survives exactly where the topology forgets.
\end{example}

\begin{boundary}
The mark mechanism is an external serialization and audit invariant. It is not a proof that a marked cover is an internal topos object, a colored-site construction in the literature, or a categorical semantics for evidence; its mathematical relationship to indexed categories and stacks has not received a dedicated novelty review. The topology consumes the unmarked family; governance code consumes the marked record. An enclosing Evidence Packet must separately bind raw ingress bytes; the cover digest is the canonical semantic-record digest only.
\end{boundary}

\section{The ambient modal candidate}\label{sec:ambient}

\begin{construction}[The successor package]\label{con:ambient}
The successor is \emph{not} defined as the already-joined sheaf category; doing so would make the ACBE and RPM reflectors identity on every admitted object and erase the modal structure a second time (Theorem~\ref{thm:resolved-identity}). Instead, begin with the ambient presheaf topos and three topology refinements:
\[
\Eamb=\PSh(\cC\times L),\qquad
J_{\mathrm{res}}=J_{\mathrm{obs}}\vee J_A\vee J_R,\qquad
\Eres=\Sh(\cC\times L,\,J_{\mathrm{res}}),
\]
where the three refinements come from legacy observational covers in the $\cC$ coordinate and the $K_A$, $K_R$ covers in the $L$ coordinate. Writing sheafification-then-inclusion as $M_{\mathrm{obs}}=i_{\mathrm{obs}}a_{\mathrm{obs}}$, $M_A=i_Aa_A$, $M_R=i_Ra_R$, and assuming $\cC$ small in the chosen universe, the underlying topos candidate and full research object are
\[
\mathrm{NoeTop}_2:=\Eamb,\qquad
\mathrm{NoeSystem}_2=\bigl(\mathrm{NoeTop}_2,\ M_{\mathrm{obs}},\ M_A,\ M_R,\ \Eres,\ \textsc{external mark record}\bigr).
\]
\end{construction}

$\mathrm{NoeSystem}_2$ is a structured package, not itself a topos or category, and the provenance record is external audit data rather than an internal topos object. The intended discipline: raw and partially resolved evidence lives in the ambient topos, acted on by the modalities; only jointly resolved objects enter $\Eres$.

\begin{boundary}
The v0.1 executable schema is deliberately narrower than Construction~\ref{con:ambient}: it represents ACBE/RPM context-factor families and their joint refinement only. Typed observational covers in the legacy $\cC$ coordinate, general arrows of $\cC\times L$, full product-site covers, and their serialization remain \textsc{open}. The notation above is a research specification, not a completed implementation.
\end{boundary}

\section{The finite-context theory}\label{sec:theorems}

Throughout this section, presheaves are ordinary small $\Set$-valued functors on $L$ in a fixed universe. All results are proved directly; the executable certificate of Section~\ref{sec:certificate} corroborates each one exhaustively.

\subsection{Subcanonicity}

\begin{theorem}[NT2-2; \textnormal{ledger \textsc{thm-nt2-subcanonical-context-factor}}]\label{thm:nt2-2}
$J_A$ and $J_R$ are subcanonical on $L$.
\end{theorem}

\begin{proof}
Let $yW$ be representable, so $yW(V)$ is a singleton when $V\subseteq W$ and empty otherwise. Fix a generating cover $K_P(V)$. A matching family assigns to each piece $V\cap B$ a section over it, compatibly. Suppose a matching family exists; then $V\cap B\subseteq W$ for every block $B$ meeting $V$. Since the blocks of $P$ partition $\Omega$, $V=\bigcup_B (V\cap B)\subseteq W$, so a unique section over $V$ exists and restricts correctly. Conversely a section over $V$ restricts to a unique matching family. The empty case is immediate (one empty family, one section of $yW(\varnothing)$, as $\varnothing\subseteq W$ always). Thus the gluing map is a bijection for every generator, and by the basis criterion (Remark~\ref{rem:background}) every representable is a sheaf.
\end{proof}

\begin{boundary}
Theorem~\ref{thm:nt2-2} concerns the context factor $L$ only. It neither repairs nor bears on the legacy subcanonicity gap of Proposition~\ref{prop:countermodel}.
\end{boundary}

\subsection{Distinctness}

\begin{definition}[Witness presheaves]\label{def:witness}
Define $F_A:L^{\mathrm{op}}\to\Set$ by: $F_A(V)$ is the set of functions from $\{i \mid A_i\subseteq V\}$ to $\{0,1\}$---bit assignments to the \emph{complete} ACBE blocks contained in $V$---with restriction maps forgetting coordinates whose blocks are no longer contained. Define $F_R$ symmetrically with the roles of $A$ and $R$ exchanged.
\end{definition}

\begin{theorem}[NT2-3: distinct, non-identity modalities; \textnormal{ledger \textsc{thm-nt2-distinct-modalities}}]\label{thm:nt2-3}
$F_A$ is a $J_A$-sheaf but not a $J_R$-sheaf, and $F_R$ is a $J_R$-sheaf but not a $J_A$-sheaf. Consequently the reflective subcategories $\Sh(L,J_A)$ and $\Sh(L,J_R)$ of $\PSh(L,\Set)$ are distinct, and neither sheafification modality is trivial.
\end{theorem}

\begin{proof}
\emph{$F_A$ satisfies ACBE gluing.} Fix $V$ and the cover $K_A(V)$. A complete block $A_i\subseteq V$ satisfies $A_i\cap A_{i'}=\varnothing$ for $i'\neq i$ and $A_i\cap A_i=A_i$, so each complete block of $V$ is a complete block of exactly one cover piece, and no block is complete in the overlap of distinct pieces (their intersection is empty, where $F_A(\varnothing)$ is the one-element set of functions from the empty index set). A matching family therefore assigns bits independently on the blocks of each piece, and these assemble uniquely to a bit assignment on the complete blocks of $V$. The gluing map is a bijection.

\emph{$F_A$ fails RPM gluing at $\Omega$.} $F_A(\Omega)$ has $2^2=4$ sections. Each RPM half $R_j$ contains no complete ACBE block ($A_i\not\subseteq R_j$ since $A_i\cap R_j$ is a singleton by \eqref{eq:transverse}), so $F_A(R_j)$ is a singleton, as is $F_A(\varnothing)$. The matching-family set for $K_R(\Omega)$ is therefore a singleton, and the gluing map $4\to 1$ is not injective. Hence $F_A$ is not a $J_R$-sheaf.

The symmetric argument applies to $F_R$. Distinctness of the subcategories follows since each contains an object the other excludes; nontriviality of $M_A$ follows because $M_A F_R\not\cong F_R$ (a reflector fixes exactly the objects of its subcategory, and $F_R\notin\Sh(L,J_A)$), and symmetrically for $M_R$.
\end{proof}

\begin{remark}[Functoriality is checked, not assumed]
The finite implementation validates functoriality of a proposed witness (all $81$ restriction squares) before checking sheaf conditions, so a nonfunctorial section assignment cannot pass as a witness. The certificate records \texttt{acbe\_witness\_functorial}, \texttt{acbe\_witness\_is\_acbe\_sheaf} true, \texttt{acbe\_witness\_is\_rpm\_sheaf} false, and the symmetric RPM facts.
\end{remark}

\begin{remark}[Ambient lift caveat]\label{rem:ambientlift}
The v0.1 specification proposed lifting distinctness to $\PSh(\cC\times L)$ by external product with a representable $y\cC_0$. That witness is false in general: at the null cover over $(c,\varnothing)$, the product $y\cC_0(c)\times F_A(\varnothing)$ need not be a singleton, and duplicated representable coordinates can obstruct gluing. A terminal $\cC$-presheaf is a candidate repair only after the lifted topology is defined, $\cC$ is assumed nonempty, and naturality is checked in both coordinates. Ambient distinctness therefore remains \textsc{open}; the proved result is scoped to $\PSh(L,\Set)$. (The ledger statement for \textsc{thm-nt2-distinct-modalities} includes the ambient-lift sentence; this paper follows the publication-layer correction and does not assert it.)
\end{remark}

\subsection{The joined topology}

\begin{lemma}[Join identification]\label{lem:join}
$J_A\vee J_R=J_{\mathrm{atom}}$, where $J_{\mathrm{atom}}$ is generated by the singleton partition $Q$ with the null cover at the empty context.
\end{lemma}

\begin{proof}
Order topologies by $J\le J'$ iff every $J$-covering sieve is $J'$-covering.

($\le$) Every singleton $\{\omega\}\subseteq V$ lies inside exactly one ACBE piece $V\cap A_i$ and one RPM piece $V\cap R_j$. Hence the sieve generated by $K_A(V)$ contains every arrow factoring through an ACBE piece---in particular every singleton inclusion---so the singleton family's sieve contains it; concretely both $K_A(V)$ and $K_R(V)$ are covers for $J_{\mathrm{atom}}$, since the sieve each generates contains the singleton family of $V$ (each piece is covered by its own singletons, and local character closes up). The null covers agree. Thus $J_A\vee J_R\le J_{\mathrm{atom}}$.

($\ge$) In any topology containing both $J_A$ and $J_R$, cover $V$ by its ACBE pieces, then cover each nonempty piece $V\cap A_i$ by its RPM pieces. By \eqref{eq:transverse} the nonempty members of the composite family are
\[
V\cap A_i\cap R_j = V\cap\{(i,j)\},
\]
the singletons of $V$; empty members do not change the generated sieve, an empty ACBE piece is refined by the null cover exactly as required, and at $V=\varnothing$ the composite is the null cover. Closure under cover composition (J3) makes the singleton-generated sieve covering. Thus $J_{\mathrm{atom}}\le J_A\vee J_R$.
\end{proof}

\begin{theorem}[NT2-4; \textnormal{ledger \textsc{thm-nt2-join-discrete-context}}]\label{thm:nt2-4}
The common refinement of the ACBE and RPM global covers is the four-singleton cover, and the joined context coverage $J_A\vee J_R$ is the topology generated by the singleton partition---equivalently, the ordinary union topology on the finite powerset lattice.
\end{theorem}

\begin{proof}
The refinement computation is \eqref{eq:transverse}; the topology identification is Lemma~\ref{lem:join}. For the final phrase: a sieve on $V$ is $J_{\mathrm{atom}}$-covering iff it contains every singleton of $V$, which is the union coverage on $\mathcal{P}(\Omega)$.
\end{proof}

\begin{boundary}
The internal label $J_{\mathrm{atom}}$ does not assert the standard category-theoretic \emph{atomic topology}; $J_{\mathrm{singleton}}$ is the preferred publication-layer name. The exact-topology implementation enumerates all $298$ sieves over the sixteen contexts, saturates maximality, pullback stability, and local character to a fixed point, verifies the Grothendieck axioms independently on both computed results, and reports $168$ covering sieves for both the join and singleton presentations.
\end{boundary}

\subsection{Explicit reflectors}

The composition theorem requires knowing sheafification exactly, not just abstractly. The candidate constructs it directly.

\begin{construction}[One-partition reflector]\label{con:reflector}
For a partition $P\in\{A,R,Q\}$ and presheaf $F$, let $I_P(V)=\{B\in P \mid V\cap B\neq\varnothing\}$ and define
\begin{equation}\label{eq:reflector}
(R_PF)(V) \;=\; \prod_{B\in I_P(V)} F(V\cap B),
\end{equation}
with the empty product the singleton $\mathbf{1}$; in particular $(R_PF)(\varnothing)=\mathbf{1}$ even when $F(\varnothing)$ has zero, two, or more elements. For $U\subseteq V$, restriction drops coordinates whose intersection with $U$ is empty and applies $F(V\cap B)\to F(U\cap B)$ on the rest; identity and composition laws hold coordinatewise. The unit is
\begin{equation}\label{eq:unit}
\eta_{P,F,V}:F(V)\to (R_PF)(V),\qquad x\mapsto\bigl(x|_{V\cap B}\bigr)_{B\in I_P(V)} .
\end{equation}
\end{construction}

\begin{lemma}\label{lem:sheafness}
$R_PF$ is a $J_P$-sheaf.
\end{lemma}

\begin{proof}
At the empty context, the null-cover sheaf condition requires exactly one section, which holds by $(R_PF)(\varnothing)=\mathbf{1}$. Let $V\neq\varnothing$ and consider $K_P(V)$. Partition blocks are disjoint, so distinct pieces meet at the empty context, where $R_PF$ has one section: compatibility on every pairwise overlap is automatic. On a nonempty piece $V\cap B$, formula \eqref{eq:reflector} has exactly one factor and is canonically $F(V\cap B)$; on an empty piece it is $\mathbf{1}$. The matching-family set is therefore exactly $\prod_{B\in I_P(V)}F(V\cap B)$, and the gluing map from $(R_PF)(V)$ is the canonical reindexing bijection. Checking basis families suffices for the generated topology.
\end{proof}

\begin{lemma}[Universal property]\label{lem:universal}
The unit \eqref{eq:unit} exhibits $R_P$ as the associated-sheaf reflector for $J_P$: for every $J_P$-sheaf $G$, composition with $\eta_{P,F}$ is a bijection
\[
\Hom_{\Sh}(R_PF,G)\;\cong\;\Hom_{\PSh}(F,G),
\]
natural in $F$ and $G$.
\end{lemma}

\begin{proof}
Given $\alpha:F\to G$, define $\bar\alpha_V$ on a tuple $(x_B)\in(R_PF)(V)$ by applying $\alpha_{V\cap B}$ componentwise; because $G(\varnothing)=\mathbf{1}$, the results form a matching family on $K_P(V)$, and the sheaf condition of $G$ glues them to a unique section $\bar\alpha_V((x_B))\in G(V)$. Restriction of a glued section is the unique gluing of the restricted family, so $\bar\alpha$ is natural; by construction $\bar\alpha\circ\eta_{P,F}=\alpha$.

Uniqueness: restrict any candidate factorization $\beta$ to a nonempty piece $V\cap B$, where the unit is the canonical singleton-factor bijection $F(V\cap B)\to(R_PF)(V\cap B)$; this forces $\beta$ to agree with $\alpha$ on every local component, and uniqueness of gluing in $G(V)$ forces $\beta_V=\bar\alpha_V$. At the empty context both sides are singletons. Naturality in $F$ and $G$ is a direct check. This direct proof avoids assuming that a single plus construction sheafifies an arbitrary presheaf; general background is \cite[Thm.~7.10.10, Prop.~7.10.12]{stacks00W1}.
\end{proof}

By Lemmas~\ref{lem:sheafness} and \ref{lem:universal}, $R_A$, $R_R$, $R_Q$ realize $M_A$, $M_R$, $M_{\mathrm{atom}}$.

\subsection{The composition theorem}

\begin{theorem}[Modality composition; \textnormal{ledger \textsc{thm-nt2-modality-composition} v0.2}]\label{thm:composition}
On $\PSh(L,\Set)$ there are unit-compatible canonical natural isomorphisms
\[
M_AM_R \;\cong\; M_{\mathrm{atom}} \;\cong\; M_RM_A ,
\]
and with Lemma~\textup{\ref{lem:join}}, either composite is naturally isomorphic to sheafification for $J_A\vee J_R$.
\end{theorem}

\begin{proof}
For every $F$ and $V$, expand \eqref{eq:reflector} twice and apply transversality \eqref{eq:transverse}:
\[
(R_AR_RF)(V)
=\!\!\prod_{\substack{i \\ V\cap A_i\neq\varnothing}}\ \prod_{\substack{j \\ V\cap A_i\cap R_j\neq\varnothing}}\!\! F(V\cap A_i\cap R_j)
\;\cong\;\prod_{\omega\in V}F(\{\omega\})
=(R_QF)(V).
\]
The middle identification flattens a nested tuple and indexes the result by the atoms of $V$; its inverse groups atomic coordinates first by $A_i$ and then by $R_j$. These are mutually inverse bijections. At the empty context both sides are the empty product $\mathbf{1}$, so the null cover requires no exception.

\emph{Naturality in $V$}: for $U\subseteq V$, both sides drop atoms not in $U$ and apply $F(\{\omega\})\to F(\{\omega\})$ identically on retained atoms; flattening commutes with restriction. \emph{Naturality in $F$}: a presheaf map $\tau:F\to F'$ acts by $\tau_{\{\omega\}}$ independently on each atomic coordinate along both routes. \emph{Unit compatibility}: the flattened composite unit sends $x\in F(V)$ to $(x|_{\{\omega\}})_{\omega\in V}$, which is exactly the unit of $R_Q$.

Swapping $A$ and $R$ gives $R_RR_A\cong R_Q$; composing one comparison with the inverse of the other yields $R_AR_R\cong R_RR_A$. By Lemmas~\ref{lem:universal} and \ref{lem:join}, $R_Q$ is sheafification for $J_A\vee J_R$. These are natural isomorphisms, not strict equalities: nested and flat $\Set$-products need not be literally equal in a chosen implementation.
\end{proof}

\begin{corollary}\label{cor:commute}
The finite-context modalities commute up to canonical natural isomorphism, and either composite lands in $\Sh(L,J_A)\cap\Sh(L,J_R)$.
\end{corollary}

\begin{corollary}\label{cor:atomicfix}
On an already atomic sheaf, both modalities and both composites are naturally isomorphic to the identity. This does not make either modality the identity on the ambient presheaf category (Theorem~\ref{thm:nt2-3}).
\end{corollary}

\begin{remark}[Distinct but commuting]\label{rem:narrative}
Theorem~\ref{thm:composition} corrects a live narrative risk in the program's earlier informal materials, which hoped for interacting, order-sensitive modal structure. The finite modalities are distinct and non-identity, yet they commute. Distinctness does not imply noncommutation; no current artifact licenses a noncommutation claim, and any successor construction seeking genuinely order-sensitive modalities must change the site again rather than reinterpret this one.
\end{remark}

\begin{theorem}[Identity on the resolved subcategory; \textnormal{ledger \textsc{thm-nt2-resolved-modalities-identity}}]\label{thm:resolved-identity}
Every sheaf for $J_A\vee J_R$ is separately a $J_A$-sheaf and a $J_R$-sheaf; consequently both reflectors restrict, up to their canonical reflector isomorphisms, to the identity on the resolved subcategory.
\end{theorem}

\begin{proof}
A sheaf for a finer topology satisfies the gluing condition for every covering sieve of any coarser topology contained in it; $J_A,J_R\le J_A\vee J_R$ by construction. A reflector is naturally isomorphic to the identity exactly on its subcategory of local objects.
\end{proof}

This theorem is the formal content of the design rule in Construction~\ref{con:ambient}: the nontrivial modal action belongs to the ambient category. No text may describe $M_A$ or $M_R$ as nontrivial endomodalities of the resolved subtopos.

\section{The executable finite certificate}\label{sec:certificate}

The finite theory above is small enough to check exhaustively, and the candidate does so with two independent executable layers.

\subsection{Finite site checker}
A dependency-free checker (\texttt{src/noetic\_topos2/finite\_site.py}, certificate \texttt{FINITE\_MODEL\_CERTIFICATE\_v0.1.json}, model ID \textsc{nt2-context-2x2-transverse-partitions}) verifies over all $4$ atoms and $16$ contexts: proper global covers and joint exhaustiveness for both partitions; pullback stability of every generating cover; sheafhood of all $16$ representables (Theorem~\ref{thm:nt2-2}); functoriality and the four sheaf/non-sheaf verdicts for both witness presheaves (Theorem~\ref{thm:nt2-3}); the transverse refinement identity (Theorem~\ref{thm:nt2-4}); coincidence of unmarked singleton families alongside distinctness of their marked records (Example~\ref{ex:singleton}); marked-cover round-trip stability, parent-provenance binding for pullbacks and joint refinements, and resolution of derived records; and both legacy collapse mechanisms (Proposition~\ref{prop:collapse}). All 23 registered checks pass.

\subsection{Modality composition verifier}
The composition layer (\texttt{modality\_composition/verify\_modality\_composition.py}) implements formula \eqref{eq:reflector}, its unit, and both flattening/regrouping maps, and computes exact finite Grothendieck topologies by enumerating every sieve on every context and saturating maximality, pullback stability, and local character to a fixed point. The deterministic run checks: all $4$ atoms, $16$ contexts, and $81$ inclusion arrows; all $298$ sieves, with $168$ covering for both the joined and singleton presentations; the Grothendieck axioms independently on both computed topologies; both index-level composite decompositions at all contexts; both canonical natural isomorphisms on a constant two-section presheaf whose empty-context value is deliberately nonterminal (exercising the $(R_PF)(\varnothing)=\mathbf{1}$ clause); a bit-assignment presheaf and all $16$ representables; and inverse laws, restriction naturality, commutation naturality, and unit compatibility across $18$ executable fixtures and $3{,}012$ restriction squares.

\subsection{Formal layer and its limits}
At the bound commit, the finite-site suite (8 tests), the modality-composition suite (6 tests), and the pinned Lean Gate~1 verifier pass. The Lean scope covers four-atom exhaustiveness, proper global pieces, transverse refinement, and witness-map lemmas. It does \emph{not} formalize the sixteen presheaves, sheaf conditions, topology generation, reflector universal properties, the composition theorem, the product-site comparison, or the legacy round trip.

\begin{boundary}
Exhaustive execution over one sixteen-object site corroborates the finite proofs; it is not proof-assistant verification, external peer review, empirical validation, or a product-claim license. This sentence is quoted from the certificate's own claim boundary and governs every number above.
\end{boundary}

\section{Gluing boundary for probabilistic and linear data}\label{sec:gluing}

The finite theory is a theory of $\Set$-valued descent, and two adjacent temptations must be blocked explicitly.

\emph{Probability.} For $\Set$-valued presheaves, gluing means bijectivity of the canonical map from global sections to matching local sections. Probability distributions do not inherit this by analogy: pairwise-compatible marginals may lack a global joint distribution (the marginal problem; nonextendable compatible families are classical \cite{vorobev}). A probabilistic successor must pass one of four explicit gates: an acyclic or junction-tree cover with a proved extension theorem; a finite linear-program feasibility certificate for the joint; an obstruction calculation with a stated coefficient system \cite{abramskybrandenburger}; or a quantitative relaxation with reported residual and uncertainty. No passing $\Set$-valued sheaf test licenses a probabilistic consistency claim.

\emph{Linear and cellular data.} A linear or cellular sheaf requires a separate linear representation---finite-dimensional stalks, chosen inner products, and weighted coboundary data---before a sheaf Laplacian is defined \cite{hansenghrist}. Nothing in this paper creates that structure, and bounded Hilbert-valued cellular sheaves need not be finite-dimensional.

\section{Backward compatibility as a conjectural comparison}\label{sec:compat}

\begin{conjecture}[Product-site round trip; \textnormal{ledger \textsc{thm-nt2-product-site-backward-compatibility}, \textsc{proof required}}]\label{conj:roundtrip}
Let $\Eleg=\Sh(\cC,\Jleg)$. After atomic resolution of the context factor there is an equivalence $\Eres\simeq\Eleg^{4}$ under which the embedding $\iota:X\mapsto(X,1,1,1)$ into a registry-fixed neutral channel $\omega_{00}$ is fully faithful and evaluation at the neutral channel satisfies a specified natural isomorphism $\mathrm{ev}_{00}\circ\iota\cong\mathrm{Id}_{\Eleg}$.
\end{conjecture}

Two distinct round trips are required and neither is complete: the categorical comparison above, and the serialization requirement that registered identifiers and canonical payload hashes survive migration by strict equality. A categorical natural isomorphism is not a bit-for-bit serialization receipt. Bare legacy site objects are sheafified before marking whenever $\Jleg$ is not known subcanonical (Proposition~\ref{prop:countermodel}). This paper claims no backward compatibility.

\begin{remark}[The expressivity boundary]\label{rem:expressivity}
Even if Conjecture~\ref{conj:roundtrip} is proved, $\Eleg^4$ is four \emph{independent} legacy channels: the discrete resolved context factor contains no remaining overlap, transition law, or cross-atom coupling---Theorem~\ref{thm:composition} lands every composite in the atomic subcategory, where Corollary~\ref{cor:atomicfix} freezes further modal action. The sixteen-context construction is therefore a finite separation and provenance reference model, not a substantively rich final semantics for contextual cognition. Registered obligation \textsc{obl-nt2-cross-context-semantics} requires a canonical choice: either (a) the ambient modal package is named the principal semantics, or (b) the resolved semantics gains typed overlap or transition data with proved coherence. Its acceptance criteria demand that the selected contract carry typed definitions, that any claimed coherence theorem have a schema-valid ledger record and reviewable proof artifact, and that the four-independent-channel limitation remain explicit unless coupling is proved. v0.1 licenses neither branch.
\end{remark}

\section{Related work and novelty status}\label{sec:related}

The site-to-topology passage, sheafification and its exactness, representable sheaves, and the Grothendieck-topos background are standard \cite{maclane,stacks00ZC,stacks00W1,stacks00WO}; no novelty is claimed for them. Sheaf-theoretic treatments of contextuality and marginal-extension obstructions \cite{abramskybrandenburger} and cellular-sheaf spectral theory \cite{hansenghrist} motivate the explicit gates of Section~\ref{sec:gluing} without being incorporated. The specific contributions here---the collapse diagnosis applied to the registered legacy clauses, the transverse two-partition finite context factor with its exact reflector calculus, and the parallel hash-chained provenance layer---are Project Coherence internal candidates whose relationship to colored sites, indexed categories, and stacks has not received a dedicated novelty search. All eight theorem-ledger records carry \textsc{novelty unassessed}; this paper accordingly does not call them new theorems in the scholarly sense.

\section{Open obligations and admission gates}\label{sec:gates}

Noetic Topos 2.0 may leave Ring~4 only after all of the following close:
\begin{enumerate}
\item independent category-theory review of Propositions~\ref{prop:collapse}--\ref{prop:countermodel} and Theorems~\ref{thm:nt2-1}--\ref{thm:composition};
\item proof-assistant formalization of the finite site, generated topologies, witness presheaves, reflector universal properties, join, and natural isomorphisms;
\item a typed construction of observational covers and general arrows in $\cC\times L$;
\item the product-site lift of Theorem~\ref{thm:composition}, its coherence in the $\cC$ coordinate, and its interaction with $J_{\mathrm{obs}}$;
\item Conjecture~\ref{conj:roundtrip} with both categorical and strict serialization round trips and tests;
\item a canonical decision under \textsc{obl-nt2-cross-context-semantics} (Remark~\ref{rem:expressivity});
\item separate probability-gluing theorems or feasibility certificates (Section~\ref{sec:gluing});
\item a dedicated novelty search for the marked/provenance construction;
\item reconciliation with the shared symbol and theorem registries; authorship, disclosure, citation, and IP review; and
\item external referee review, with explicit confirmation that no empirical or product wording was licensed.
\end{enumerate}

\section{Conclusion}

The conservative correction and the successor construction fit together as one argument. The legacy ACBE/RPM clauses collapse because identities enter their sieves (Proposition~\ref{prop:collapse}); no relabeling escapes this, so the successor changes the generators to proper transverse partition pieces on a new finite context factor. There, the theory is complete and exact: genuine subcanonical topologies (Theorems~\ref{thm:nt2-1}--\ref{thm:nt2-2}), distinct non-identity modalities witnessed by explicit presheaves (Theorem~\ref{thm:nt2-3}), a computed join (Theorem~\ref{thm:nt2-4}), explicit reflectors verified against the universal property (Lemmas~\ref{lem:sheafness}--\ref{lem:universal}), and a commuting composition law (Theorem~\ref{thm:composition})---each exhaustively corroborated by the finite certificate. Provenance that the topology necessarily forgets survives in hash-chained marked records that fail closed. What remains open is exactly what the ledger says is open: the ambient lift, the legacy comparison, richer cross-context semantics, probabilistic descent, formalization, novelty assessment, and independent review. Preserving those boundaries with the same precision as the theorems is part of the result.

\begin{thebibliography}{9}
\bibitem{maclane} S.~Mac Lane and I.~Moerdijk, \emph{Sheaves in Geometry and Logic: A First Introduction to Topos Theory}, Springer, 1994.
\bibitem{stacks00ZC} The Stacks Project Authors, ``Sites and sheaves: The topology defined by a site,'' Tag 00ZC, \url{https://stacks.math.columbia.edu/tag/00ZC}.
\bibitem{stacks00W1} The Stacks Project Authors, ``Sheafification,'' Tag 00W1 (esp.\ Theorem 7.10.10, Proposition 7.10.12), \url{https://stacks.math.columbia.edu/tag/00W1}.
\bibitem{stacks00WO} The Stacks Project Authors, ``Representable sheaves,'' Tag 00WO, \url{https://stacks.math.columbia.edu/tag/00WO}.
\bibitem{vorobev} N.~N. Vorob'ev, ``Consistent families of measures and their extensions,'' \emph{Theory Probab.\ Appl.}\ 7 (1962), 147--163.
\bibitem{abramskybrandenburger} S.~Abramsky and A.~Brandenburger, ``The sheaf-theoretic structure of non-locality and contextuality,'' \emph{New J.\ Phys.}\ 13 (2011), 113036.
\bibitem{hansenghrist} J.~Hansen and R.~Ghrist, ``Toward a spectral theory of cellular sheaves,'' \emph{J.\ Appl.\ Comput.\ Topol.}\ 3 (2019), 315--358.
\end{thebibliography}

\end{document}
